% ---------------------------------------------------------------------------
% howtoread.tex
% ---------------------------------------------------------------------------
\chapter*{How to Read This Booklet}
\addcontentsline{toc}{chapter}{How to Read This Booklet}
\markboth{How to Read This Booklet}{How to Read This Booklet}

\section*{Typographic conventions}

\begin{itemize}
  \item \code{Monospaced text} is code: an identifier, a keyword, a file name
        or a command.
  \item \file{Green monospaced text} is a file in the accompanying source
        tree.
  \item \tool{Sans-serif} names are tools: \tool{QuestaSim}, \tool{Verilator},
        \tool{cocotb}, \tool{GTKWave}.
\end{itemize}

Listings that carry line numbers are complete, reproduced in full from
\file{code/cordic/}; nothing here is excerpted.

\section*{How to read this one}

Read it with the sources open in an editor, not on their own. \Cref{ch:case}
is written as one continuous build: specification, fixed-point format, the
VHDL implementation, the SystemVerilog implementation, then testbench after
testbench. It assumes you already have the vocabulary from
\refdesignbook, \refverificationbook and \refsimulationbook; where a
construct is used here without re-explanation, that is where it was
introduced.

If you only want to see one testbench style, the chapter is written so that
each testbench section is self-contained once the design itself is
understood: you do not have to read all six to follow one.

\section*{The companion booklets}

This is the \textbf{Example Cookbook}. \refdesignbook and
\refverificationbook cover the language used here; \refsimulationbook covers
the tools (\tool{Verilator}, \tool{QuestaSim}, \tool{cocotb}) that run the
testbenches built here.
